% TEMPLATE-MILC-v3.tex - plantilla maestra UNICA de las guias MILC v3.
%
% La usan las cuatro familias de guias, cada una con su driver:
%   · Grados 6-11          -> scripts/build-guias-g11.py
%   · Bebras               -> scripts/build-guias-bebras.py
%   · Semillero            -> scripts/build-guias-semillero.py
%   · Territorio interior  -> scripts/build-guias-territorio-interior.py
%
% Antes habia tres copias casi identicas (~400 lineas iguales, 6 distintas):
% cada mejora habia que aplicarla tres veces, y en la practica no se hacia
% (el motor de recursos vivia solo en dos de ellas). Lo que varia entre
% familias esta parametrizado en 6 placeholders que cada driver llena
% (se nombran aqui SIN su sintaxis real: el builder reemplaza por texto plano,
% tambien dentro de los comentarios, y un valor multilinea rompe el preambulo):
%
%   HEADER_IZQ        encabezado izquierdo de pagina
%   HEADER_DER        encabezado derecho de pagina
%   PORTADA_ETIQUETA  rotulo sobre el numero grande (GRADO/MOMENTO/MODULO)
%   PORTADA_SERIE     linea de serie de la portada
%   PORTADA_ID        identificador de la guia en la portada
%   PORTADA_CIRCULO   el circulito de grado (°), o vacio si no aplica
%   PDF_TITULO        titulo en texto plano (sin \textbf ni comillas TeX) para
%                     los metadatos del PDF (/Title); lo produce lib_xelatex.texto_pdf
%
% Composicion (auditoria 2026-09-03): las bandas de estacion piden espacio
% (\Needspace*) y prohiben el salto justo despues (\nopagebreak), asi no quedan
% huerfanas al pie; las cajas largas (softbox, drawbox, infoband, puentebox) son
% `breakable`, asi una caja de media pagina ya no arrastra todo a la siguiente
% dejando la anterior al 40 %. Todo texto sobre mostaza o verde va en milcNegro
% (blanco sobre #E5B400 da 1,93:1 y sobre #58A923 2,95:1; ninguno pasa WCAG AA).
% El resto de claves las documenta content/guias/_SCHEMA.md. El script valida
% que no queden placeholders sin reemplazar antes de escribir el archivo final.

% Etiquetado PDF (latex-lab): /Lang, estructura y texto alternativo de las
% figuras, para lectores de pantalla. Debe ir ANTES de \documentclass.
\DocumentMetadata{lang=es-CO,tagging=on}
\documentclass[11pt,letterpaper]{article}
% headheight=24pt: la cabecera derecha lleva dos lineas (identidad de la guia
% y estacion actual). headsep se acorta en la misma medida para que el bloque
% de texto quede exactamente donde estaba con headheight=12pt.
\usepackage[margin=2.0cm,headheight=24pt,headsep=13pt]{geometry}
\usepackage[table]{xcolor}
\usepackage{fontspec}
\usepackage[spanish]{babel}
\usepackage{tikz}
% Librerías TikZ para los diagramas de las guías (bloque `recursos:`):
%  · babel  → babel en español vuelve activos caracteres como «>», y TikZ los usa
%             en sus opciones (>=stealth, ->). Sin esta librería, cualquier
%             diagrama con flechas revienta con «\language@active@arg> has an
%             extra }». Debe ir ANTES de que se dibuje nada.
%  · positioning        → sintaxis «below left=2pt and 3pt».
%  · decorations.pathreplacing → llaves ({brace}) para señalar rangos.
\usetikzlibrary{babel,positioning,decorations.pathreplacing}
\usepackage{graphicx}
% Circuitos: el trabajo de electrónica se hace hoy con micro:bit y sus sensores
% integrados (MakeCode, sin cableado), así que no se necesitan esquemáticos.
% Si algún día entran montajes con protoboard, basta descomentar esta línea:
% los diagramas .tex/.tikz declarados en `recursos:` ya se incrustan solos.
% \usepackage{circuitikz}
\usepackage[most]{tcolorbox}
\usepackage{tabularx}
\usepackage{array}
\usepackage{fancyhdr}
\usepackage{titlesec}
\usepackage[document]{ragged2e}  % bandera a la derecha en todo el cuerpo
\usepackage{enumitem}
\usepackage{microtype}
\usepackage{etoolbox}
\usepackage{needspace}
% hyperref al final del preambulo: metadatos del PDF (titulo, autor, idioma,
% familia), marcadores por estacion (\pdfbookmark) y enlaces sin recuadro.
\usepackage[hidelinks,
  pdftitle={SUMA, PROMEDIO, MAX y MIN — cuatro cifras para decidir, nunca una sola},
  pdfauthor={PhD. Álvaro Cárdenas Orozco},
  pdfsubject={Tecnología e Informática · Grado 8 · Guía 4 · MILC v3},
  pdflang={es-CO},
  bookmarksopen=true,bookmarksnumbered=false]{hyperref}

% Helvetica Neue no trae la flecha U+2192, asi que cada «→» escrito literalmente
% en el YAML se compilaba a NADA, en silencio (462 flechas en 61 guias: rutas del
% tipo «paleta → tipografia → lockup» salian rotas). Mapeamos el caracter a su
% version matematica, que si tiene glifo. Asi el autor puede seguir escribiendo
% «→» comodamente en el YAML.
\usepackage{newunicodechar}
\newunicodechar{→}{\ensuremath{\rightarrow}}
% Mismo problema con los simbolos de marca y vinetas: el «✓» de «una palomita ✓»
% salia como cuadrito vacio en el PDF de todas las guias que lo usaban.
\usepackage{pifont}
\newunicodechar{✓}{\ding{51}}
\newunicodechar{✗}{\ding{55}}
\newunicodechar{✔}{\ding{52}}
\newunicodechar{•}{\textbullet}
\newunicodechar{≥}{\ensuremath{\geq}}
\newunicodechar{≤}{\ensuremath{\leq}}

\setmainfont{Helvetica Neue}
\setsansfont{Helvetica Neue}
\setlength{\parindent}{0pt}
\setlength{\parskip}{6pt}
% Sin particion de palabras: en 6.º-7.º una palabra cortada por guion al final
% de linea («Comuni-cacion») frena la lectura mucho mas que una linea corta.
\hyphenpenalty=10000
\exhyphenpenalty=10000
\pretolerance=2000
\tolerance=3000
\setlength{\emergencystretch}{3em}
\linespread{1.10}
\renewcommand{\arraystretch}{1.25}
\setlist{leftmargin=5mm,itemsep=1mm,topsep=1mm}
\newcolumntype{Y}{>{\RaggedRight\arraybackslash}X}

\definecolor{milcMagenta}{HTML}{E6007E}
\definecolor{milcVino}{HTML}{5A0038}
\definecolor{milcTurquesa}{HTML}{0093A5}
\definecolor{milcVerde}{HTML}{58A923}
\definecolor{milcMostaza}{HTML}{E5B400}
\definecolor{milcOcre}{HTML}{B86600}
\definecolor{milcNegro}{HTML}{241816}
\definecolor{milcGris}{HTML}{F7F7F4}
\definecolor{milcLinea}{HTML}{E8E2DD}
\definecolor{milcGradePrimary}{HTML}{FF6600}
\definecolor{milcGradeDark}{HTML}{9A3E00}
\definecolor{milcGradeSoft}{HTML}{FFE4D0}
\definecolor{milcGradeAccent}{HTML}{CFE7FF}
\definecolor{milcGradeText}{HTML}{FFFFFF}
\color{milcNegro}

\pagestyle{fancy}
\fancyhf{}
% Cabecera de dos lineas: identidad de la guia arriba y, a la derecha y en
% gris, la estacion en curso (\markright desde \milcestacion). Antes de la
% primera estacion la segunda linea va vacia; el \strut mantiene la altura
% para que la linea de arriba no baile entre paginas.
\lhead{\small Tecnología e Informática · Grado 8\\[-1pt]{\footnotesize\strut}}
\rhead{\small Guía 4 · MILC v3\\[-1pt]{\footnotesize\strut\color{milcNegro!65}\rightmark}}
\cfoot{\small\thepage}
\renewcommand{\headrulewidth}{0pt}

% Sobre mostaza (#E5B400) y verde (#58A923) el texto va en milcNegro; sobre el
% resto de la paleta, en blanco. Se usa dentro de fonttitle/fontupper, donde un
% \color posterior manda sobre coltitle.
\newcommand{\milctintasobre}[1]{%
  \ifboolexpr{test{\ifstrequal{#1}{milcMostaza}} or test{\ifstrequal{#1}{milcVerde}}}%
    {\color{milcNegro}}{\color{white}}}

\newtcolorbox{titlebox}[1]{
  enhanced,colback=#1,colframe=#1,arc=7mm,boxrule=0pt,
  left=7mm,right=7mm,top=3mm,bottom=3mm,
  before skip=8pt,after skip=8pt,
  fontupper=\sffamily\bfseries\Large\color{white}
}

\newtcolorbox{softbox}[3]{
  enhanced,breakable,pad at break=3mm,
  colback=#2,colframe=#1,borderline west={4pt}{0pt}{#1},
  arc=2mm,boxrule=.4pt,left=4mm,right=4mm,top=3mm,bottom=3mm,
  before skip=6pt,after skip=8pt,title={#3},coltitle=white,
  colbacktitle=#1,fonttitle=\sffamily\bfseries\milctintasobre{#1},fontupper=\RaggedRight,
  attach boxed title to top left={xshift=3mm,yshift=-2mm},
  boxed title style={arc=2mm, boxrule=0pt}
}

\newtcolorbox{drawbox}[2]{
  enhanced,breakable,pad at break=4mm,
  colback=white,colframe=#1,arc=4mm,boxrule=1pt,
  left=5mm,right=5mm,top=4mm,bottom=4mm,before skip=8pt,after skip=8pt,
  title={#2},coltitle=white,colbacktitle=#1,fonttitle=\sffamily\bfseries\milctintasobre{#1},
  fontupper=\RaggedRight,
  attach boxed title to top left={xshift=4mm,yshift=-2mm},
  boxed title style={arc=3mm, boxrule=0pt}
}

\newtcolorbox{infoband}[1]{
  enhanced,breakable,pad at break=2mm,colback=#1!8,colframe=#1!50,arc=2mm,boxrule=.5pt,
  left=4mm,right=4mm,top=2mm,bottom=2mm,before skip=4pt,after skip=4pt,
  fontupper=\small\RaggedRight
}

% Puente narrativo entre secciones (contrato editorial Conéctate).
% Da continuidad: "Acabas de X. Ahora vas a Y porque Z."
% Visualmente: banda izquierda en color del grado, texto en itálica pequeña.
\newtcolorbox{puentebox}{
  enhanced,breakable,pad at break=2mm,colback=milcGradeSoft,colframe=milcGradeDark,
  borderline west={3pt}{0pt}{milcGradeDark},
  arc=0mm,boxrule=0pt,left=5mm,right=4mm,top=2.5mm,bottom=2.5mm,
  before skip=4pt,after skip=8pt,
  fontupper=\itshape\small\color{milcNegro}\RaggedRight
}
% La flecha va en modo matematico a proposito: Helvetica Neue NO tiene el glifo
% U+2192, asi que \textrightarrow se compilaba a NADA («Missing character: There
% is no → in font Helvetica Neue») y el puente salia sin flecha. En modo
% matematico la flecha viene de la fuente matematica y si se dibuja.
\newcommand{\puente}[1]{%
  \begin{puentebox}%
    \textcolor{milcGradeDark}{$\rightarrow$}\hspace{2mm}#1%
  \end{puentebox}%
}

\newcommand{\linead}[1]{\noindent\color{milcLinea}\rule{#1}{0.5pt}\color{milcNegro}}
\newcommand{\respuesta}[1]{\par\vspace{2mm}\foreach \n in {1,...,#1}{\noindent\color{milcLinea}\rule{\linewidth}{0.5pt}\par\vspace{5mm}}\color{milcNegro}}
\newcommand{\checkbox}{\textcolor{milcGradeDark}{\fbox{\phantom{x}}}\hspace{5pt}}

% ─── Listas aireadas para las guías socioemocionales ────────────────────
% Componer las verificaciones como «\checkbox texto\\» deja el texto que salta
% de línea debajo del cuadrito y apelmaza el bloque. Estos entornos alinean la
% continuación y airean los ítems. Los usa Territorio Interior; cualquier
% familia puede hacerlo.
\newlist{checklista}{itemize}{1}
\setlist[checklista]{
  label=\textcolor{milcGradeDark}{\fbox{\phantom{x}}},
  leftmargin=9mm, labelsep=3.2mm, itemsep=2.4mm,
  topsep=2.5mm, parsep=0pt, partopsep=0pt, before=\RaggedRight
}
\newlist{pasoslista}{enumerate}{1}
\setlist[pasoslista]{
  label=\textcolor{milcGradeDark}{\sffamily\bfseries\arabic*.},
  leftmargin=8mm, labelsep=3mm, itemsep=2.8mm,
  topsep=1mm, parsep=0pt, partopsep=0pt, before=\RaggedRight
}

\newcommand{\phasecard}[4]{
\begin{tcolorbox}[enhanced,colback=#3,colframe=#2,arc=3mm,boxrule=.5pt,height=3.05cm,valign=center,halign=center,left=1.5mm,right=1.5mm,top=2mm,bottom=2mm]
{\sffamily\bfseries\fontsize{22}{24}\selectfont\textcolor{#2}{#1}}\par
{\sffamily\bfseries\small #4}
\end{tcolorbox}
}

% ── Piezas del rediseno 2026 (legibilidad 6.º-11.º) ───────────────────
% Banda de estacion: reemplaza el titlebox macizo. Titulo a la izquierda,
% dato practico (tiempo/modalidad) a la derecha, en una sola linea.
% Nunca huerfana: pide 9 lineas libres (o salta de pagina rellenando con
% \vfill) y prohibe el salto justo despues de la banda. Nueve y no seis:
% la caja partible que sigue necesita ~80pt (titulo adosado, relleno y dos
% lineas) para arrancar en la misma pagina; con menos, tcolorbox la manda
% entera a la siguiente y la banda se queda sola. El penalty 10000 va
% en `after`, ANTES del espacio posterior: un `after skip` (aunque sea 0pt)
% es un \vskip tras la caja, y ese glue es un punto de corte legal que un
% \nopagebreak escrito despues de \end{tcolorbox} ya no cierra. Cada banda
% deja marcador PDF y actualiza la cabecera derecha.
\newcounter{milcestacion}
\newcommand{\milcestacion}[2]{%
\Needspace*{9\baselineskip}%
\stepcounter{milcestacion}%
\pdfbookmark[1]{#2}{milcestacion\themilcestacion}%
\markright{#2}%
\begin{tcolorbox}[enhanced,colback=#1,colframe=#1,arc=2mm,boxrule=0pt,
  left=5mm,right=5mm,top=2.6mm,bottom=2.6mm,before skip=11pt,
  after={\par\nopagebreak\vskip7pt}]
{\sffamily\bfseries\fontsize{15}{18}\selectfont\milctintasobre{#1}#2}
\end{tcolorbox}}

% Tarjeta de paso numerada: sustituye las tablas «Pilar / Aplicacion», que
% eran formato de consulta docente, no de estudiante. El numero va en un
% circulo sobre el borde para que la secuencia se lea de un vistazo.
\newcommand{\milcpaso}[3]{%
\Needspace{4\baselineskip}%
\begin{tcolorbox}[enhanced,colback=white,colframe=milcLinea,arc=1.5mm,boxrule=.9pt,
  left=14mm,right=4mm,top=3mm,bottom=3mm,before skip=4pt,after skip=4pt,
  fontupper=\RaggedRight,
  overlay={\node[anchor=north west,circle,fill=#2,text=white,inner sep=0pt,
    minimum size=7.4mm,font=\sffamily\bfseries\small]
    at ([xshift=3.6mm,yshift=-3.2mm]frame.north west) {#1};}]
#3
\end{tcolorbox}}

% Casilla marcable de tamano util con lapiz (la anterior era \fbox{\phantom{x}},
% de alto variable segun el texto que la seguia).
\newcommand{\milcmarca}{\framebox[3.8mm]{\rule{0pt}{3.4mm}}}

% Fila marcable de lista suelta (checks de la estacion 1).
\newcommand{\milccheckrow}[1]{%
\par\vspace{1.8mm}\noindent
\makebox[6.4mm][l]{\raisebox{-.55mm}{\textcolor{milcNegro}{\milcmarca}}}%
\parbox[t]{\dimexpr\linewidth-6.4mm\relax}{\RaggedRight #1}\par}

% Celda marcable dentro de la rubrica (tabla de autoevaluacion). Misma
% sangria francesa, pero pensada para vivir dentro de una columna X.
\newcommand{\milcopcion}[1]{%
\makebox[5.2mm][l]{\raisebox{-.55mm}{\textcolor{milcNegro}{\milcmarca}}}%
\parbox[t]{\dimexpr\linewidth-5.2mm\relax}{\RaggedRight #1}}

% ── Recursos visuales de la guía ──────────────────────────────────────
% Declarados en el bloque `recursos:` del YAML y emitidos por el builder.
% \guiaFigura   → imagen rasterizada o PDF (foto, carta celeste, montaje).
% \guiaDiagrama → fuente TikZ/circuitikz (.tex) incrustada como vector:
%                 es la vía para circuitos y esquemáticos editables.
% [#1] = texto alternativo (alt) · #2 = ruta relativa al .tex · #3 = pie
% El alt llega al PDF etiquetado (/Alt) y lo lee el lector de pantalla.
\newcommand{\guiaFigura}[3][]{%
\begin{center}
\includegraphics[width=0.88\linewidth,height=0.38\textheight,keepaspectratio,alt={#1}]{#2}\\[1.5mm]
{\footnotesize\itshape\textcolor{milcNegro!75}{#3}}
\end{center}
}

\newcommand{\guiaDiagrama}[2]{%
\begin{center}
\input{#1}\\[1.5mm]
{\footnotesize\itshape\textcolor{milcNegro!75}{#2}}
\end{center}
}

\begin{document}
\thispagestyle{empty}

% ── Portada ───────────────────────────────────────────────────────────
% Antes: un rectangulo de color a pagina completa. Se veia bien en pantalla,
% pero en la fotocopiadora del colegio se comia el toner y salia gris sucio.
% Ahora la banda ocupa el tercio superior y el resto va en blanco.
\begin{tikzpicture}[remember picture,overlay]
  \path[fill=milcGradePrimary]
    (current page.north west) rectangle ([yshift=-10.6cm]current page.north east);

  \node[anchor=north west,text=milcGradeText] at ([xshift=2.0cm,yshift=-1.55cm]current page.north west)
    {\sffamily\bfseries\fontsize{12}{14}\selectfont GRADO};
  \node[anchor=north west,text=milcGradeText,opacity=.85] at ([xshift=2.0cm,yshift=-2.35cm]current page.north west)
    {\sffamily\bfseries\fontsize{8.5}{10}\selectfont SERIE GUÍAS · TECNOLOGÍA E INFORMÁTICA};
  \node[anchor=north east,text=milcGradeText,opacity=.85,align=right] at ([xshift=-2.0cm,yshift=-1.55cm]current page.north east)
    {\sffamily\bfseries\fontsize{9}{12}\selectfont I.E. Sor María Juliana\\Cartago, Valle del Cauca};

  \node[anchor=west,text=milcGradeText] (gradonum) at ([xshift=1.85cm,yshift=-7.1cm]current page.north west)
    {\sffamily\bfseries\fontsize{112}{112}\selectfont 8};
  % El circulito de grado (°) solo aplica a las guias de grado («11°»); en
  % Bebras y Semillero el numero grande es un momento/modulo, no un grado.
  % Se posiciona relativo al nodo `gradonum`, no en coordenadas absolutas.
  \draw[milcGradeText,line width=1.6pt]
    ([xshift=2.0mm,yshift=-4.5mm]gradonum.north east) circle (.15cm);

  \node[anchor=west,text=milcGradeText,text width=11.4cm,align=left] at ([xshift=1.5cm,yshift=0.5cm]gradonum.east)
    {
     {\sffamily\bfseries\fontsize{9}{11}\selectfont\MakeUppercase{Período 1 · Análisis de datos con phronesis}}\\[3mm]
     {\sffamily\bfseries\fontsize{21}{25}\selectfont\setlength{\baselineskip}{27pt}Cuatro cifras\\[4pt]para decidir}\\[3.5mm]
     {\sffamily\bfseries\fontsize{15}{18}\selectfont Guía 4-8-TIC}
    };
\end{tikzpicture}

\vspace*{9.4cm}
\vfill

{\sffamily\bfseries\fontsize{8}{10}\selectfont\textcolor{milcGradeDark}{LO QUE TE LLEVAS HOY}}

\begin{tcolorbox}[enhanced,colback=white,colframe=milcNegro,arc=2mm,boxrule=1.6pt,
  left=5mm,right=5mm,top=4mm,bottom=4mm,before skip=3pt,after skip=12pt,
  fontupper=\RaggedRight\fontsize{11}{15.5}\selectfont]
Una hoja de Excel con las cuatro funciones (SUMA, PROMEDIO, MAX, MIN) aplicadas a una columna real de tu tabla limpia. Al lado de cada cifra, una decisión escrita en dos líneas. Y un párrafo de cinco líneas sobre cuál cifra sirvió más para decidir.
\end{tcolorbox}

\begin{tcolorbox}[enhanced,colback=milcGris,colframe=milcGris,arc=2mm,boxrule=0pt,
  left=5mm,right=5mm,top=3.5mm,bottom=3.5mm,after skip=12pt]
{\sffamily\bfseries\fontsize{8}{10}\selectfont\textcolor{milcNegro!55}{ESTA GUÍA ES DE}}\\[3mm]
\begin{tabularx}{\linewidth}{X p{3.2cm} p{3.2cm}}
\linead{\linewidth} & \linead{3.0cm} & \linead{3.0cm}\\[-1mm]
{\fontsize{8.5}{10}\selectfont\textcolor{milcNegro!55}{Nombre completo}} &
{\fontsize{8.5}{10}\selectfont\textcolor{milcNegro!55}{Grupo}} &
{\fontsize{8.5}{10}\selectfont\textcolor{milcNegro!55}{Fecha}}\\
\end{tabularx}
\end{tcolorbox}

\vfill

\noindent\textcolor{milcLinea}{\rule{\linewidth}{1pt}}\\[2mm]
{\sffamily\bfseries\fontsize{9.5}{12}\selectfont Autor: Dr. Álvaro Cárdenas Orozco}\\[1mm]
{\sffamily\fontsize{8.5}{11}\selectfont\textcolor{milcNegro!55}{Metodología MILC · apertura ancestral · pensamiento computacional · triángulo Dussel–estoicismo–Floridi}}

\newpage

\section*{Apertura: SUMA, PROMEDIO, MAX y MIN --- cuatro cifras para decidir, nunca una sola}

\begin{softbox}{milcGradeDark}{milcGradeSoft}{Saber ancestral}
El río La Vieja bordea a Cartago y le pone frontera con Risaralda. Cuando crece, no se lleva la ciudad entera. Se lleva siempre los mismos barrios: Brisas del Río, La Platanera, La Arenera, La Playa, Tierra del Olvido. En marzo de 2022 quedaron bajo el agua cuando el río llegó a 12,3 metros (CiudadRegión, 2022). La gente de la orilla no adivina la creciente: la mira. Puerto Alejandría es el punto al que se va a ver el nivel. Allí se compara cómo está el río hoy contra cómo estaba ayer, y contra la marca de la última vez que se salió. Esa marca es un \textbf{máximo}: el nivel más alto que se recuerda. Y hay también un \textbf{mínimo}: el nivel de verano, cuando se ve la arena. Ninguno de los dos se decide con el promedio. La \textbf{cara de exclusión} es dura: la creciente golpea por clase y por ubicación, no por azar, y quien vive lejos del río no ve la marca. Hoy vas a calcular cuatro cifras sobre tu tabla. Y vas a aprender que el promedio, solo, engaña como engañaría mirar el río «en promedio» cuando ya está a 12 metros.\par\smallskip{\footnotesize\textcolor{milcNegro!70}{\textbf{Fuente.} CiudadRegión. (2022, 7 de marzo). \emph{Alarma por alto nivel del río La Vieja a su paso por Cartago}. CiudadRegión Noticias.}}
\end{softbox}

\begin{softbox}{milcTurquesa}{milcGris}{Saber contemporáneo}
Excel tiene cuatro funciones que resumen una columna entera en una sola cifra. \texttt{=SUMA(rango)} responde «¿cuánto en total?». \texttt{=PROMEDIO(rango)} responde «¿cuánto es lo típico?». \texttt{=MAX(rango)} responde «¿cuál fue el caso más alto?». \texttt{=MIN(rango)} responde «¿cuál fue el más bajo?». Se llaman \textbf{medidas resumen}: en vez de mirar 200 filas, miras una cifra que las representa. El error más común es mirar solo el promedio. Un promedio de notas de 3,5 puede venir de un grupo parejo o de un grupo donde alguien sacó 5,0 y alguien sacó 1,0. La cifra es la misma; la situación es distinta. Por eso el oficio pide \textbf{las cuatro juntas}: cada una responde una pregunta que las otras tres no responden.
\end{softbox}

\begin{softbox}{milcMostaza}{milcGris}{Pregunta-puente}
\textit{En Puerto Alejandría nadie mira el nivel promedio del río: mira el máximo de la última creciente y el nivel de hoy. Cuando calcules el promedio de las notas del salón, ¿qué te va a esconder ese número? ¿Y qué te dirían el máximo y el mínimo que el promedio no dice?}
\end{softbox}

\begin{softbox}{milcVerde}{milcGris}{Saber y saber hacer}
Al terminar podrás: (1) \textbf{identificar} qué pregunta responde cada una de las cuatro funciones; (2) \textbf{aplicar} las cuatro en Excel con la sintaxis correcta sobre tu tabla limpia; (3) \textbf{analizar} qué dice cada cifra y qué esconde; (4) \textbf{evaluar} cuál de las cuatro sirvió más para tomar una decisión concreta.
\end{softbox}



\small
\begin{tabularx}{\textwidth}{>{\bfseries}p{3.0cm}Y}
\rowcolor{milcGradePrimary!12}Autor & Dr. Álvaro Cárdenas Orozco\\
DBA & Tratamiento de datos cuantitativos con criterio ético (MEN, T\&I)\\
Referentes & Phronesis aristotélica · Floridi (ética del dato) · Estoicismo (ver lo que es)\\
Producto final & Una hoja de Excel con las cuatro funciones (SUMA, PROMEDIO, MAX, MIN) aplicadas a una columna real de tu tabla limpia. Al lado de cada cifra, una decisión escrita en dos líneas. Y un párrafo de cinco líneas sobre cuál cifra sirvió más para decidir.\\
\end{tabularx}
\normalsize

\puente{En la sesión 3 dejaste una tabla limpia con bitácora. Hoy la haces hablar con cuatro cifras. En la sesión 5 vas a copiar fórmulas sin romperlas, y para eso necesitas haber escrito hoy tus primeras funciones.}

\milcestacion{milcGradeDark}{Ruta MILC: así vamos a aprender}
\begin{tabularx}{\textwidth}{>{\centering\arraybackslash}X>{\centering\arraybackslash}X>{\centering\arraybackslash}X>{\centering\arraybackslash}X}
\phasecard{1}{milcVerde}{milcGris}{Escucha\\[-1mm]\small Observo, escucho y pregunto.} &
\phasecard{2}{milcTurquesa}{milcGris}{Entender\\[-1mm]\small Sistematizo y ordeno.} &
\phasecard{3}{milcMagenta}{milcGris}{Hacer\\[-1mm]\small Construyo y aplico.} &
\phasecard{4}{milcVino}{milcGris}{Cierre\\[-1mm]\small Evalúo y reflexiono.}
\end{tabularx}


\puente{Antes de que Excel calcule por ti, vas a calcular tú, a mano, con diez valores reales. Así sabes qué hace cada función antes de que la máquina lo haga en un segundo.}

\milcestacion{milcVerde}{Estación 1 de 3 · Escucha}

\begin{softbox}{milcVerde}{milcGris}{Escena para escuchar}
\textbf{Actividad 1 · IDENTIFICA --- Las cuatro cuentas a mano} (15 min · individual). (1) Toma diez valores reales de tu tabla limpia: diez notas, diez gastos, diez temperaturas o los minutos que dormiste en diez días. (2) Con calculadora, sin Excel, calcula la suma de los diez. (3) Divide esa suma entre diez: ese es el promedio. (4) Busca el valor más alto y el más bajo: son el máximo y el mínimo. (5) Escribe una línea: ¿cuál de las cuatro cifras te dice más sobre tu situación?
\end{softbox}

{\sffamily\bfseries\fontsize{8}{10}\selectfont\textcolor{milcNegro!55}{MARCA CUANDO LO HAYAS HECHO}}
\milccheckrow{Tengo los diez valores y las cuatro cifras calculadas a mano.}
\milccheckrow{La suma es mayor que el máximo y el promedio está entre el mínimo y el máximo.}
\milccheckrow{Escribí cuál de las cuatro cifras me dice más y por qué.}

\begin{infoband}{milcVerde}


\textbf{Lo que va al cuaderno (Actividad 1 · IDENTIFICA).} \textbf{Título:} «Actividad 1 --- Las cuatro cuentas a mano». \textbf{Formato:} ficha con los diez valores, las cuatro cifras y una línea sobre cuál dice más. \textbf{Extensión:} media página. \textbf{Sabes que terminaste cuando} las cuatro cifras están, la suma es mayor que el máximo, y puedes decir qué hace cada función sin mirar la guía.
\end{infoband}

\puente{Ya sentiste las cuatro cuentas en la calculadora. Ahora las escribes en Excel sobre tu tabla limpia, con tu pareja, y aprendes la trampa del promedio solo.}

\milcestacion{milcTurquesa}{Estación 2 de 3 · Entender}

\begin{softbox}{milcTurquesa}{milcGris}{Lo que vas a entender}
\textbf{Actividad 2 · APLICA --- Las cuatro funciones en Excel} (25 min · parejas). Ya sabes qué hacen. Ahora que las haga Excel. (1) Con tu pareja, abran la tabla limpia de la sesión 3. (2) Elijan una columna numérica: notas, gastos, ventas, temperaturas. (3) En cuatro celdas debajo de la columna escriban \texttt{=SUMA(}, \texttt{=PROMEDIO(}, \texttt{=MAX(} y \texttt{=MIN(}, seleccionen el rango y cierren el paréntesis. (4) En la celda de la izquierda de cada una escriban qué es: «Total», «Promedio», «Máximo», «Mínimo». (5) Comprueben que la suma es mayor que el máximo y que el promedio está entre el mínimo y el máximo.
\end{softbox}

\milcpaso{1}{milcTurquesa}{\textbf{Las cuatro funciones.} \texttt{=SUMA(B2:B21)} suma los números del rango: ¿cuánto en total? \texttt{=PROMEDIO(B2:B21)} suma y divide entre la cantidad de valores: ¿cuánto es lo típico? \texttt{=MAX(B2:B21)} devuelve el más alto: ¿cuál fue el mejor caso? \texttt{=MIN(B2:B21)} devuelve el más bajo: ¿cuál fue el peor?}
\milcpaso{2}{milcTurquesa}{\textbf{La trampa del promedio solo.} Dos salones con promedio 3,5. En uno todos sacaron entre 3,0 y 4,0. En el otro hay un 5,0 y un 1,0. El promedio no distingue. El máximo y el mínimo sí. Por eso nunca se mira una cifra sola.}
\milcpaso{3}{milcTurquesa}{\textbf{Lo que el rango esconde.} Si el rango incluye el encabezado, la función lo ignora, pero si incluye una celda con texto, también la ignora sin avisar. Y el promedio no cuenta las celdas vacías, pero sí cuenta los ceros. Lo que limpiaste en la sesión 3 decide lo que sale hoy.}
\milcpaso{4}{milcTurquesa}{\textbf{La pregunta decide la función.} «¿Cuánto recaudó la tienda esta semana?» es SUMA. «¿Qué tan típica fue la semana?» es PROMEDIO. «¿Cuál fue el día más fuerte?» es MAX. «¿Cuál fue el más flojo?» es MIN. Primero la pregunta; después la función.}

\begin{drawbox}{milcTurquesa}{Las cuatro funciones, en una mirada}
\textbf{SUMA:} cuánto en total. \\[1mm] \textbf{PROMEDIO:} cuánto es lo típico. \\[1mm] \textbf{MAX:} el caso más alto. \\[1mm] \textbf{MIN:} el caso más bajo. \\[1mm] \textbf{Regla:} las cuatro juntas, nunca una sola.
\end{drawbox}

\begin{softbox}{milcMostaza}{milcGris}{Errores comunes y su corrección}
(1) \textbf{Mirar solo el promedio}: esconde los extremos. (2) \textbf{Rango mal seleccionado}: dejar fuera la última fila o meter una celda de otra columna. (3) \textbf{Texto dentro del rango}: la función lo salta sin avisar y la suma sale menor. (4) \textbf{Confundir promedio con lo que le pasa a la mayoría}: un solo valor extremo mueve el promedio. (5) \textbf{Cifra sin etiqueta}: un 35,5 suelto en una celda no significa nada.
\end{softbox}

\begin{infoband}{milcTurquesa}


\textbf{Lo que va al cuaderno (Actividad 2 · APLICA).} \textbf{Título:} «Actividad 2 --- Las cuatro funciones en Excel». \textbf{Formato:} tabla de 4 filas y 3 columnas (función / fórmula escrita / resultado). \textbf{Extensión:} media página. \textbf{Sabes que terminaste cuando} las cuatro fórmulas están escritas con su rango, cada resultado tiene etiqueta en Excel, y la comprobación de coherencia da bien.
\end{infoband}

\puente{Ya tienes las cuatro cifras. Toca decidir: al lado de cada una escribes qué dice, qué sugiere y qué harías tú. Lo escribes tú, con tus datos.}

\milcestacion{milcMagenta}{Estación 3 de 3 · Hacer}

\begin{softbox}{milcMagenta}{milcGris}{Cómo vas a construir tu producto}
\textbf{Actividad 3 · EVALÚA --- Una decisión por cada cifra} (30 min · individual). Cuatro números no son una decisión. (1) Vuelve a tu hoja y mira las cuatro cifras. (2) Al lado de cada una escribe dos líneas: qué dice el dato y qué harías tú con eso. (3) Compara el máximo y el mínimo con el promedio: ¿el grupo es parejo o disparejo? Escríbelo. (4) Cierra con un párrafo de cinco líneas: cuál de las cuatro cifras sirvió más para decidir, y por qué. (5) Guarda la hoja como G8-P1-S4-funciones.xlsx.
\end{softbox}

\milcpaso{1}{milcMagenta}{\textbf{Tu entrega tiene tres piezas.} (1) Las cuatro funciones con etiqueta en la hoja. (2) Cuatro decisiones de dos líneas, una por cifra. (3) El párrafo de cinco líneas sobre cuál cifra sirvió más.}
\milcpaso{2}{milcMagenta}{\textbf{Del dato a la acción.} «La SUMA de gastos fue 50.000 pesos esta semana. Eso dice que gasté más de lo que tenía. Yo cambiaría dos días de tienda por almuerzo de casa.» Tres frases: dato, lectura, acción. Sin la tercera, el número queda huérfano.}
\milcpaso{3}{milcMagenta}{\textbf{Parejo o disparejo.} Si el máximo y el mínimo están cerca del promedio, el grupo es parejo y el promedio lo representa bien. Si están lejos, el promedio esconde a alguien. Y ese alguien es quien más necesita que lo miren.}
\milcpaso{4}{milcMagenta}{\textbf{Antes de entregar, léelo en voz alta.} Si una decisión suena a «hay que mejorar», no es una decisión: es un deseo. Una decisión dice quién hace qué, cuándo. Reescríbela hasta que diga eso.}

\begin{drawbox}{milcMagenta}{Antes de entregar}
\checkbox Las cuatro funciones escritas con el rango correcto. \\[1mm] \checkbox Cada resultado con etiqueta a la izquierda. \\[1mm] \checkbox La suma es mayor que el máximo; el promedio está entre el mínimo y el máximo. \\[1mm] \checkbox Cuatro decisiones de dos líneas, con dato, lectura y acción. \\[1mm] \checkbox Párrafo de cinco líneas sobre la cifra que sirvió más. \\[1mm] \checkbox Hoja guardada como G8-P1-S4-funciones.xlsx.
\end{drawbox}

\begin{softbox}{milcMagenta}{milcGris}{Plantilla de trabajo}
\textbf{SUMA.} \textit{Resultado, qué dice, qué haría.} \\[1mm] \textbf{PROMEDIO.} \textit{Resultado, qué dice, qué haría.} \\[1mm] \textbf{MAX.} \textit{Resultado, qué dice, qué haría.} \\[1mm] \textbf{MIN.} \textit{Resultado, qué dice, qué haría.} \\[1mm] \textbf{Cierre.} \textit{La cifra que sirvió más fue… porque…}
\end{softbox}

\begin{infoband}{milcMagenta}


\textbf{Lo que va al cuaderno (Actividad 3 · EVALÚA).} \textbf{Título:} «Actividad 3 --- Una decisión por cada cifra». \textbf{Formato:} tabla de 4 filas y 3 columnas (cifra / qué dice / qué haría), más el párrafo de cierre. \textbf{Extensión:} una página. \textbf{Sabes que terminaste cuando} cada cifra tiene su acción, dijiste si el grupo es parejo o disparejo, y el párrafo nombra la cifra que sirvió más.
\end{infoband}

\puente{Lo que entregas hoy: las cuatro funciones sobre una columna real, cuatro decisiones escritas y un párrafo sobre cuál cifra sirvió más. La prueba de calidad: ninguna cifra queda sola, sin decisión al lado.}

\milcestacion{milcMostaza}{Producto de la sesión}
\textbf{Cuatro funciones con etiqueta + cuatro decisiones + párrafo sobre la cifra que sirvió más}.
\begin{softbox}{milcMostaza}{milcGris}{Criterios de calidad}
(1) Las cuatro funciones están escritas con el rango correcto y con etiqueta. (2) La suma es mayor que el máximo y el promedio está entre el mínimo y el máximo. (3) Cada cifra tiene una decisión de dos líneas con dato, lectura y acción. (4) Dijiste si el grupo es parejo o disparejo comparando máximo, mínimo y promedio. (5) El párrafo de cierre nombra la cifra que sirvió más y por qué. (6) La hoja está guardada con nombre claro.
\end{softbox}

\puente{Antes de cerrar, mira lo que hiciste desde cinco lados. Una cifra sin decisión es un número huérfano. Una decisión sin cifra es una corazonada.}

\milcestacion{milcVino}{Evaluación liberadora · cinco dimensiones}

\begin{softbox}{milcVino}{milcGris}{Sentido de la evaluación}
Evaluar no es solo revisar una nota. Es reconocer qué aprendí, qué cambió en mí, qué emociones manejé, cómo me relaciono con la ciudad y la comunidad, y qué le dejaré a quien viene detrás.
\end{softbox}

\textbf{Mi reflexión liberadora en cinco dimensiones.} Completa cada frase iniciada:

\small
\begin{tabularx}{\textwidth}{>{\bfseries\RaggedRight\arraybackslash}p{3.4cm}Y>{\RaggedRight\arraybackslash}p{4.6cm}}
\rowcolor{milcVino}\color{white}Dimensión & \color{white}\bfseries Mi respuesta (frame starter) & \color{white}\bfseries Algo que descubrí\\
1. Personal & Lo que más cambió en mí fue \linead{6.5cm} & \linead{4.2cm}\\
2. Emocional & La emoción más fuerte fue \linead{2.5cm} y la regulé \linead{3cm} & \linead{4.2cm}\\
3. Ciudadana & En democracia esto importa porque \linead{6cm} & \linead{4.2cm}\\
4. Local (Cartago) & En mi barrio sería útil para \linead{6.5cm} & \linead{4.2cm}\\
5. Intergeneracional & A mi abuela/abuelo le diría \linead{6.5cm} & \linead{4.2cm}\\
\end{tabularx}
\normalsize

\begin{infoband}{milcVino}
\textbf{Para la próxima sustentación.} Vuelve a esta tabla en seis meses. Verás que las respuestas cambian — no porque lo aprendido cambie, sino porque tú cambias. La evaluación liberadora es una conversación contigo a través del tiempo.
\end{infoband}


\puente{Para terminar, tres ideas sobre mirar el conjunto, contadas con palabras de esta guía: Dussel sobre el que queda fuera del promedio, Marco Aurelio sobre mirar el conjunto y no solo un punto, Floridi sobre los patrones pequeños que solo valen si se comparan.}

\milcestacion{milcGradeDark}{Tres ideas para cerrar · Dussel, un estoico y Floridi}

\begin{softbox}{milcVino}{milcGris}{Enrique Dussel · lente del nosotros}
\textit{El que queda fuera del sistema tiene un rostro que reclama. Mirarlo es el comienzo de cualquier justicia.}

{\footnotesize\itshape\color{milcNegro!70}--- idea tomada de Enrique Dussel, contada con palabras de esta guía. No es una cita textual.}

El promedio del salón no tiene rostro. El mínimo sí: es alguien concreto que sacó la nota más baja. Un análisis que solo mira el promedio deja a esa persona fuera del sistema, aunque esté en la tabla.

\textbf{¿Quién es el mínimo de mi tabla, y qué decisión lo tiene en cuenta?}
\end{softbox}
\linead{\linewidth}\\[1mm]
\linead{\linewidth}

\begin{softbox}{milcMostaza}{milcGris}{Estoicismo · Marco Aurelio · cuidado interior}
\textit{Mira el curso completo de las cosas, como quien gira con los astros, y no te quedes en un solo punto.}

{\footnotesize\itshape\color{milcNegro!70}--- idea tomada de Marco Aurelio, contada con palabras de esta guía. No es una cita textual.}

Una cifra sola es un punto. Las cuatro juntas son el curso completo: cuánto hay, cuánto es lo típico, hasta dónde sube y hasta dónde baja. Decidir con una sola es decidir a ciegas.

\textbf{¿Estoy mirando las cuatro cifras, o me quedé con la que me convenía?}
\end{softbox}
\linead{\linewidth}\\[1mm]
\linead{\linewidth}

\begin{softbox}{milcTurquesa}{milcGris}{Luciano Floridi · lente de la infoesfera}
\textit{Un patrón pequeño solo significa algo si se agregó bien, se comparó con otro y llegó a tiempo.}

{\footnotesize\itshape\color{milcNegro!70}--- idea tomada de Luciano Floridi, contada con palabras de esta guía. No es una cita textual.}

Tu promedio es un patrón pequeño. Solo vale si el rango estaba limpio (agregado bien), si lo comparaste con el máximo y el mínimo, y si decides con él esta semana y no dentro de un año.

\textbf{¿Con qué comparé mi cifra antes de decidir?}
\end{softbox}
\linead{\linewidth}\\[1mm]
\linead{\linewidth}

\begin{infoband}{milcNegro}
\textbf{Nota para el docente.} Las tres ideas están contadas con palabras de esta guía a partir del pensamiento de cada autor; no son citas textuales. De 9.º en adelante se trabajan con la obra y su referencia.
\end{infoband}

\milcestacion{milcVino}{Mi autoevaluación · marca una casilla por fila}

{\small
\setlength{\extrarowheight}{2.2pt}
\begin{tabularx}{\textwidth}{>{\RaggedRight\arraybackslash\bfseries}p{3.0cm}YYY}
\rowcolor{milcVino}\color{white}\bfseries Criterio & \color{white}\bfseries Lo logré & \color{white}\bfseries En proceso & \color{white}\bfseries Necesito apoyo\\[.6mm]
Comprensión del tema & \milcopcion{Explico las ideas con mis palabras.} & \milcopcion{Reconozco las ideas, falta precisión.} & \milcopcion{Necesito repasar lo básico.}\\[.8mm]
\rowcolor{milcGris}Pensamiento computacional & \milcopcion{Usé los pilares al armar mi producto.} & \milcopcion{Usé algunos pilares.} & \milcopcion{Necesito releer la estación 2.}\\[.8mm]
Aplicación y producto & \milcopcion{Mi producto cumple los criterios.} & \milcopcion{Está armado, con ajustes pendientes.} & \milcopcion{Aún está en construcción.}\\[.8mm]
\rowcolor{milcGris}Reflexión 5 dimensiones & \milcopcion{Respondí cada una con honestidad.} & \milcopcion{Respondí algunas con detalle.} & \milcopcion{Mis respuestas son muy generales.}\\[.8mm]
Triángulo de pensamiento & \milcopcion{Conecté las 3 voces con mi trabajo.} & \milcopcion{Conecté al menos una.} & \milcopcion{No las relaciono con mi proceso.}\\[.8mm]
\end{tabularx}}

\vspace{3mm}
\textbf{Hoy entendí que:}\\[1mm]
\linead{\linewidth}\\[1mm]
\linead{\linewidth}

\textbf{Mi compromiso para la próxima sesión es:}\\[1mm]
\linead{\linewidth}\\[1mm]
\linead{\linewidth}

\begin{softbox}{milcMostaza}{milcGris}{Cita de despedida · Marco Aurelio}
\textit{``El obstáculo en el camino se vuelve el camino. Lo que estorba al camino se vuelve camino.''}\\[1mm]
Cada error de hoy, bien observado, se convierte en pieza del camino que pisarás mañana.
\end{softbox}

\small
\begin{tabularx}{\textwidth}{>{\bfseries}p{3.6cm}Y}
\rowcolor{milcGradeDark!12}Nombre completo: & \linead{10cm}\\
Fecha: & \linead{4cm} \quad Curso/grupo: \linead{4cm}\\
Mi firma: & \linead{9cm}\\
\end{tabularx}
\normalsize

\begin{infoband}{milcGradeDark}
\textbf{Para la próxima sesión.} Llega con tu hoja de las cuatro funciones y las cuatro decisiones. En la sesión 5 vas a copiar fórmulas a muchas celdas sin que se rompan, con referencias relativas y absolutas. Las funciones de hoy son las que vas a copiar.
\end{infoband}

\end{document}
